% ── 2. Related Work ───────────────────────────────────────────────────────────
\section{Related Work}
\label{sec:related}

\paragraph{Mathematical reasoning benchmarks.}
The dominant paradigm for evaluating mathematical ability in LLMs relies on end-to-end accuracy.
GSM8K \citep{cobbe2021training} tests grade-school word problems; MATH \citep{hendrycks2021measuring} covers competition-level algebra, geometry, and number theory; MathBench \citep{liu2024mathbench} provides multi-level curricula; MathVista \citep{lu2024mathvista} extends evaluation to visual contexts; and selected BIG-Bench tasks \citep{srivastava2023beyond} probe arithmetic and logical reasoning.
Recent diagnostic work has exposed fragilities behind high aggregate scores: GSM-Symbolic \citep{mirzadeh2024gsmsymbolic} shows that adding a single irrelevant clause can drop accuracy by up to 65\%, and MATH-Perturb \citep{huang2025mathperturb} demonstrates that structural perturbations defeat memorised solution paths.
Chain-of-thought prompting \citep{wei2022chain} and its variants \citep{wang2023selfconsistency,yao2023tree} have substantially raised scores on all of these benchmarks.
However, the shared limitation is that they treat every correct answer identically regardless of the reasoning path that produced it.
A model that cross-multiplies two fractions and a model that recognises both are near $\nicefrac{1}{2}$ receive the same credit.
This conflation obscures whether LLMs possess efficient, flexible numerical reasoning or merely execute brute-force procedures.
\textsc{SenseMath} addresses this gap by controlling the \emph{availability} of shortcut strategies across matched item variants, so that accuracy differences are attributable to strategy selection rather than problem difficulty.

\paragraph{Efficient and compressed reasoning.}
A growing line of work aims to preserve reasoning quality while reducing token expenditure.
Chain of Draft \citep{kim2024chain} generates minimal intermediate steps instead of verbose chains.
Sketch-of-Thought \citep{aytes2025sketch} uses cognitive-inspired paradigms to reduce token use by up to 84\%.
Token-Budget-Aware reasoning \citep{han2024tokenbudget} inserts explicit budgets into prompts, compressing verbose reasoning by 67\% with no accuracy loss.
\citet{chen2024overthinking} show that O1-like models waste computation on trivial problems, precisely the failure mode that number-sense shortcuts are designed to avoid.
These approaches share a common assumption: the underlying \emph{strategy} is fixed (typically step-by-step computation), and the goal is to express it more concisely.
By contrast, number sense involves switching to a qualitatively different strategy (estimation, benchmarking, or structural decomposition) that is inherently shorter because it skips unnecessary computation.
Our work is therefore orthogonal: efficient-reasoning methods compress the \emph{same} strategy; \textsc{SenseMath} measures whether models can \emph{select a different} strategy.

\paragraph{Number sense in mathematics education.}
The concept of number sense originates in the mathematics education literature.
\citet{mcintosh1992proposed} define it as ``a person's general understanding of number and operations along with the ability and inclination to use this understanding in flexible ways to make mathematical judgements.''
\citet{reys1991number} identify key components including magnitude estimation, use of benchmarks (e.g., $\nicefrac{1}{2}$, powers of 10), and recognition of numerical structure (e.g., near-cancellation, compatible-number pairs).
\citet{sowder1992making} emphasise that number sense is not a single skill but a disposition to choose efficient paths. We would now call this a \emph{meta-cognitive strategy selection} ability.
\citet{yang2003assessing} operationalise these ideas into validated test instruments for school-age children.
Crucially, this literature demonstrates that expert human reasoners do not uniformly apply algorithmic procedures; they adaptively select shortcuts when the problem structure permits.
We adopt this taxonomy as the category system for \textsc{SenseMath}.

\paragraph{Numerical cognition in LLMs.}
Recent work has begun probing whether LLMs possess numerical understanding beyond surface-level pattern matching.
\citet{rahman2025fragile} test LLMs on escalating numerical tasks and find strong deterministic performance but failure on heuristic-search problems, suggesting a ``fragile number sense.''
\citet{li2025numeracy} introduce NumericBench covering six numeracy capabilities and expose systematic gaps in basic numerical skills.
At the mechanistic level, \citet{nikankin2024arithmetic} show that LLMs solve arithmetic using sparse neuron-level heuristics rather than learned algorithms, providing a circuit-level explanation for why models may or may not exploit human-like shortcuts.
\citet{yan2025grasp} find that models achieve high numeric accuracy on addition but violate commutativity in up to 20\% of cases, directly relevant to our equation-reasoning category.
\textsc{SenseMath} complements these investigations by testing whether shortcut-sensitive reasoning can be \emph{selectively activated} through prompting, rather than probing for its unconditional presence.

\paragraph{Diagnostic and contrastive evaluation.}
Standard accuracy-based benchmarks can be ``gamed'' by surface-level cues.
Contrast sets \citep{gardner2020evaluating} address this by perturbing evaluation examples to create minimally different instances that change the gold label, exposing model brittleness.
CheckList \citep{ribeiro2020beyond} provides a task-agnostic framework of behavioural tests (invariance, directional expectation, and minimum functionality) applied to NLP models.
BIG-Bench Hard \citep{suzgun2023challenging} isolates a subset of BIG-Bench tasks where LLMs initially fail, enabling targeted analysis of failure modes.
These works share the philosophy that controlled variation is more informative than aggregate accuracy.
\textsc{SenseMath} extends this philosophy to mathematical reasoning: rather than perturbing surface text, we perturb the \emph{numerical structure} of problems to toggle shortcut availability while holding surface complexity constant.
The family-matching design (strong-shortcut, weak-shortcut, and control variants sharing the same template and digit scale) directly parallels the contrast-set methodology, but targets strategy selection rather than label robustness.

\paragraph{Positioning of \textsc{SenseMath}.}
To our knowledge, \textsc{SenseMath} is the first benchmark that \emph{controls shortcut availability at the item level} within matched families.
Existing math benchmarks vary difficulty but not strategy affordance; existing contrastive benchmarks perturb text but not numerical structure; existing efficiency work compresses a fixed strategy rather than evaluating strategy choice.
\textsc{SenseMath} occupies the intersection: each item family holds surface complexity fixed while toggling whether a number-sense shortcut is viable, enabling causal attribution of performance differences to strategy selection.
The family-matching design further supports within-family statistical tests (permutation-based asymmetry tests) that are impossible with unpaired item sets.
